\documentclass[12pt]{jlreq}
\usepackage{amsmath, amssymb}

\newcommand{\C}{\mathbb{C}}
\newcommand{\R}{\mathbb{R}}
\newcommand{\Z}{\mathbb{Z}}
\newcommand{\Tr}{\operatorname{Tr}}
\newcommand{\Impart}{\operatorname{Im}}
\newcommand{\AF}{\operatorname{AF}}
\newcommand{\EG}{\operatorname{EG}}
\newcommand{\AX}{\operatorname{AX}}
\newcommand{\GL}{\operatorname{GL}}
\newcommand{\Hom}{\operatorname{Hom}}
\newcommand{\Ker}{\operatorname{Ker}}
\newcommand{\Id}{\operatorname{Id}}
\newcommand{\Sym}{\operatorname{Sym}}

\begin{document}

平面内の長方形領域 $D = \{ (x,y) \in \mathbb{R}^2 \mid 0 < x < 1,\ 0 < y < b \}$ の境界を $\Gamma$ とおく。ただし $b$ は正の数で，$b^2$ は無理数であるとする。$D$ 上で次の境界値問題を考える。
\[
  (*) \quad \begin{cases} -\Delta u = \lambda u, & (x,y) \in D \\ u = 0, & (x,y) \in \Gamma \end{cases}
\]
ここで $\lambda$ は実数の定数である。また，$\Delta$ は $D$ 上の $C^2$ 級関数 $f(x,y)$ に対して
\[
  (\Delta f)(x,y) = \frac{\partial^2 f}{\partial x^2}(x,y) + \frac{\partial^2 f}{\partial y^2}(x,y)
\]
で定められる偏微分作用素とする。実数 $\lambda$ に対して以下の $3$ つの条件 (i), (ii), (iii) を満たす実数値関数 $u$ が存在するとき，$\lambda$ は $(*)$ の固有値であるという。
\begin{enumerate}
  \item[(i)] 関数 $u$ は，$D$ 上 $C^2$ 級で，$D \cup \Gamma$ 上連続である。
  \item[(ii)] 関数 $u$ は $D$ で恒等的に $0$ ではない。
  \item[(iii)] $\lambda$ に対して関数 $u$ は $(*)$ を満たす。
\end{enumerate}
そのとき，$u$ は $(*)$ の固有関数であるという。以下の問に答えよ。

\begin{enumerate}
  \item $u(x,y) = v(x)w(y)$ という変数分離型で表される $(*)$ の固有関数，およびそれに対応する固有値をすべて求めよ。
  \item $(*)$ の固有値と固有関数は，(1) で与えたもので尽くされることを示せ。
  \item 実数 $\mu > 0$ に対し，$\lambda \le \mu$ を満たす $(*)$ の固有値の個数を $N(\mu)$ と表すことにする。$N(\mu)/\mu^\alpha$ が $\mu \to \infty$ のとき $0$ でない有限の値に収束するように実数 $\alpha$ を定めよ。また，そのときの極限値を求めよ。
\end{enumerate}

\end{document}